% problem-type: truyen-song-3d-cauchy-thuan-nhat
% order: 20
% Bo cuc: De vi du -> Cong thuc -> De + Bai giai

\section{Truyền sóng 3D thuần nhất --- Kirchhoff}

% ---------------------------------------------------------------
\subsection*{Đề bài ví dụ}
% ---------------------------------------------------------------
Giải bài toán Cauchy cho phương trình truyền sóng 3D \textbf{thuần nhất} (không nguồn)
\[
  u_{tt}=\Delta u,\quad (x,y,z)\in\mathbb{R}^3,\ t>0,\qquad
  u(\cdot,0)=0,\quad u_t(\cdot,0)=\chi_{B_1(0)}.
\]
($B_1(0)$ là hình cầu đơn vị tâm gốc.) \textbf{Tính $u(x_0,t)$ tại điểm $x_0$ cách gốc $|x_0|=2$.}
\emph{Vì dữ liệu đối xứng cầu quanh gốc nên $u$ chỉ phụ thuộc khoảng cách $|x_0|$ --- không cần biết tọa độ cụ thể của $x_0$, chỉ cần $|x_0|=2$.}

\emph{\small(Ví dụ tự soạn minh hoạ riêng công thức Kirchhoff thuần nhất; kỹ thuật này là phần ``dữ liệu ban đầu'' của các đề sóng 3D có nguồn: MAT3365 HK2 2023--24 Câu 3, MAT3365 HK1 2025--26 Câu 3.)}

% ---------------------------------------------------------------
\subsection*{Nhận ra dạng này khi nào?}
% ---------------------------------------------------------------
\begin{itemize}
  \item PDE sóng \textbf{không nguồn} ($f=0$) trên $\mathbb{R}^3$, dữ liệu ban đầu $(g,h)$.
  \item Số chiều $n=3$ (lẻ) $\Rightarrow$ dùng Kirchhoff; nghiệm có nguyên lý Huygens mạnh (không đuôi).
\end{itemize}
\begin{meobox}
Chỉ cần ráp $g,h$ vào công thức Kirchhoff. Mẹo: nghiệm chỉ ``sống'' khi mặt cầu $\partial B_t(x)$ còn chạm giá dữ liệu.
\end{meobox}

% ---------------------------------------------------------------
\subsection*{Công thức \& phương pháp}
% ---------------------------------------------------------------
\subsubsection*{Công thức Kirchhoff}
% method: cong-thuc-kirchhoff
\begin{dieukienbox}
Dùng khi: phương trình sóng \textbf{thuần nhất} ($u_{tt}=c^2\Delta u$, không nguồn) trên $\mathbb{R}^3$.
Chỉ đúng cho $n=3$ chiều không gian. Vận tốc lan truyền là $c$.
\end{dieukienbox}

\begin{nhobox}
\textbf{Dữ liệu ban đầu:} $g(x)=u(x,0)$ (vị trí lúc đầu), $\;h(x)=u_t(x,0)$ (vận tốc lúc đầu).

\textbf{Công thức Kirchhoff} ($n=3$, vận tốc $c$) --- tích phân trên mặt cầu \textbf{bán kính $ct$}:
\[
  u(x,t)=\partial_t\!\left(\frac{1}{4\pi c^2 t}\int_{\partial B_{ct}(x)} g\,dS\right)
  +\frac{1}{4\pi c^2 t}\int_{\partial B_{ct}(x)} h\,dS .
\]
Mặt cầu $\partial B_{ct}(x)$ tâm $x$ bán kính $ct$ có diện tích $4\pi(ct)^2=4\pi c^2t^2$, nên
$\dfrac{1}{4\pi c^2t^2}\int_{\partial B_{ct}(x)}\!\cdots\,dS$ chính là \textbf{giá trị trung bình} trên mặt cầu đó.

Trường hợp $g=0$ (đứng yên lúc đầu) $\Rightarrow$ chỉ còn:
\[
  u(x,t)=\frac{1}{4\pi c^2 t}\int_{\partial B_{ct}(x)} h\,dS .
\]
\end{nhobox}

\begin{luuybox}
Đừng quên $c$: bán kính mặt cầu là $ct$ (\emph{không} phải $t$), hệ số là $\dfrac{1}{4\pi c^2t}$. Khi $c=1$ mới rút về $\dfrac{1}{4\pi t}$, mặt cầu bán kính $t$. (Đề thật hay gặp $c^2=4,9$ tức $c=2,3$.)
\end{luuybox}

\begin{meobox}
Nghiệm tại $(x,t)$ \emph{chỉ} phụ thuộc dữ liệu nằm \emph{đúng trên mặt cầu} $|y-x|=ct$ (không phải bên trong).
Vậy nếu mặt cầu này không chạm vùng mà $g,h$ khác $0$ thì $u(x,t)=0$.
\end{meobox}

\emph{(Cách tính $\int_{\partial B_t}h\,dS$: xem Phụ lục C --- Kỹ thuật tính toán.)}

\subsubsection*{Nguyên lý Huygens mạnh}
% method: phan-tich-nguyen-ly-huygens
\begin{dieukienbox}
Dùng khi: $n=3$ (số chiều lẻ), phương trình sóng thuần nhất.
Giúp \emph{thu hẹp} vùng dữ liệu thực sự ảnh hưởng đến nghiệm tại $(x_0,t)$:
chỉ còn \emph{mặt cầu} $|y-x_0|=t$, không phải toàn bộ quả cầu.
\end{dieukienbox}

\begin{nhobox}
\textbf{Nguyên lý Huygens mạnh} ($n=3$):
\[
  u(x_0,t) \text{ chỉ phụ thuộc } g,h \text{ trên mặt cầu } |y-x_0|=t.
\]
Hệ quả trực tiếp: nếu mặt cầu $\partial B_t(x_0)$ không cắt $\operatorname{supp}(g)\cup\operatorname{supp}(h)$
thì $u(x_0,t)=0$.

Tín hiệu \emph{bật lên rồi tắt hẳn} --- không có đuôi sóng kéo dài.
\end{nhobox}

\begin{meobox}
So sánh nhanh: $n=2$ (chẵn) có đuôi sóng (tín hiệu lưu lại sau khi sóng đi qua).
$n=3$ (lẻ): sóng đi qua $\Rightarrow$ im lặng hoàn toàn.

\medskip
Với phần nguồn (qua Duhamel): tích phân theo $s$ khiến có thể có ``đuôi'' theo $s$ ---
phân tích riêng cho từng lát $f(\cdot,s)$.
\end{meobox}


% ---------------------------------------------------------------
\subsection*{Đề bài \& bài giải}
% ---------------------------------------------------------------
\paragraph{Nhắc lại đề.} $u_{tt}=\Delta u$, $g=0$, $h=\chi_{B_1(0)}$. Tính $u(x_0,t)$ với $|x_0|=2$.

\paragraph{Bước 1 --- Kirchhoff.} Vì $g=0$, nghiệm bằng tích phân mặt của $h$ trên mặt cầu $\partial B_t(x_0)$:
\[
  u(x_0,t)=\frac{1}{4\pi t}\int_{\partial B_t(x_0)} \chi_{B_1(0)}\,dS
        =\frac{1}{4\pi t}\cdot\big(\text{diện tích phần }\partial B_t(x_0)\text{ nằm trong }B_1(0)\big).
\]

\paragraph{Bước 2 --- tính tích phân mặt bằng chỏm cầu} (kỹ thuật ở Phụ lục C). Mặt cầu bán kính
$R=t$ tâm $x_0$ cắt quả cầu dữ liệu $B_1(0)$ (bán kính $\rho=1$, \emph{tâm} $c=0$ là gốc). Khoảng cách
hai tâm $d=|x_0-c|=|x_0-0|=|x_0|=2$ (đúng bằng khoảng cách điểm hỏi tới gốc mà đề cho):
\[
  \cos\theta=\frac{R^2+d^2-\rho^2}{2Rd}=\frac{t^2+4-1}{4t}=\frac{t^2+3}{4t}.
\]
Có giao khi $\cos\theta<1\Leftrightarrow t^2-4t+3<0\Leftrightarrow 1<t<3$. Diện tích chỏm cầu
(công thức $S=2\pi R^2(1-\cos\theta)$, Phụ lục C), rút gọn $1-\tfrac{t^2+3}{4t}=\tfrac{4t-t^2-3}{4t}=\tfrac{(t-1)(3-t)}{4t}$:
\[
  S=2\pi t^2\Big(1-\frac{t^2+3}{4t}\Big)=2\pi t^2\cdot\frac{(t-1)(3-t)}{4t}=\frac{\pi t}{2}(t-1)(3-t).
\]

\paragraph{Bước 3 --- ráp lại.} Chia diện tích chỏm cầu cho $4\pi t$:
\begin{nhobox}
\[
  u(x_0,t)=\frac{1}{4\pi t}\cdot S=\begin{cases} \dfrac{(t-1)(3-t)}{8}, & 1\le t\le 3,\\[6pt] 0, & \text{còn lại}.\end{cases}
\]
Tín hiệu đến $x_0$ tại $t=1$ (sóng vừa chạm quả cầu), đỉnh $\tfrac18$ tại $t=2$, rồi \textbf{tắt hẳn} sau
$t=3$ --- đúng nguyên lý Huygens mạnh (sóng 3D không có đuôi).
\end{nhobox}
